\section*{الفصل \ref{ch:g11:plate-motions} --- الصفائح وحركاتها}

\begin{solution}{exo:g11:plate-motions:1}
\omterm{def:g11:plate-motions:rigid-plate}{الصفيحة الصلبة} قطعة من \omterm{def:g11:lithosphere-asthenosphere:lithosphere}{الغلاف الصخري} تتحرك ككتلة واحدة: فالمسافات
بين نقاطها لا تتغير، وتشوهها مجموع في
نطاقات ضيقة على امتداد حدودها. والزلازل تبيّن ذلك: فكلها تقريبًا
تقع على امتداد أحزمة ضيقة، هي حدود الصفائح، وقليل جدًا منها داخل
الصفائح.
\end{solution}

\begin{solution}{exo:g11:plate-motions:2}
\qty{72}{mm} في السنة تساوي \qty{72}{km} في كل مليون سنة. وفي 30 سنة تتحرك المحطة
$72 \times 30 = \qty{2160}{mm}$، أي نحو \qty{2.2}{m}؛ وفي مليون سنة،
\qty{72}{km}.
\end{solution}

\begin{solution}{exo:g11:plate-motions:3}
$v_E = 30/10 = \qty{3}{cm}$ في السنة و$v_N = 40/10 = \qty{4}{cm}$ في السنة.
والسرعة $\sqrt{3^2 + 4^2} = \qty{5}{cm}$ في السنة. والزاوية $\beta$ مع
الشمال تحقق $\tan\beta = 3/4$، إذن $\beta \approx \ang{37}$: فالمحطة
تتحرك نحو الشمال الشرقي، في اتجاه يميل \ang{37} شرق الشمال.
\end{solution}

\begin{solution}{exo:g11:plate-motions:4}
\omterm{def:g11:plate-motions:hotspot-track}{أثر البقعة الساخنة} هو خط البراكين والجبال البحرية الذي تحمله صفيحة
بعيدًا عن بقعة ساخنة وهي تتحرك فوقها، وهي تزداد قِدمًا كلما ابتعدت عن
البركان النشط. وأحدث بركان، كيلاويا، يقع عند الطرف الجنوبي
الشرقي للسلسلة، في جزيرة هاواي؛ وتزداد الأعمار
نحو الغرب الشمالي الغربي. والصفيحة تحمل براكينها بعيدًا عن
البقعة الساخنة في اتجاه حركتها: فالصفيحة الهادئية تتحرك نحو
الغرب الشمالي الغربي.
\end{solution}

\begin{solution}{exo:g11:plate-motions:5}
$519/5.1 \approx \qty{102}{km}$ في كل مليون سنة، أي \qty{102}{mm} في
السنة، أو نحو \qty{10}{cm} في السنة.
\end{solution}

\begin{solution}{exo:g11:plate-motions:6}
$v_E = 63/31 \approx \qty{20.3}{mm}$ في السنة و$v_N = 48/31 \approx
\qty{15.5}{mm}$ في السنة. والسرعة $\sqrt{20.3^2 + 15.5^2} \approx
\qty{25.5}{mm}$ في السنة؛ و$\tan\alpha = 15.5/20.3 \approx 0.76$ يعطي $\alpha
\approx \ang{37}$: فالمحطة تتحرك في اتجاه يميل \ang{37} شمال الشرق، نحو
الشمال الشرقي. وماونا كيا تتحرك أسرع بنحو $72/25.5 \approx 2.8$ مرة.
\end{solution}

\begin{solution}{exo:g11:plate-motions:7}
(1) بالملليمترات في السنة،
\[
  \vec v_{\mathrm{NA/EU}} = (-10.5 - 7.7,\ 20.6 - 16.3) = (-18.2, +4.3).
\] ومقدارها $\sqrt{18.2^2 + 4.3^2} \approx
\qty{18.7}{mm}$ في السنة، و$\tan\alpha = 4.3/18.2 \approx 0.24$ يعطي
اتجاهًا يميل \ang{13} شمال الغرب. (2) تتباعد الصفيحتان عبر الظهرة بنحو
\qty{1.9}{cm} في السنة: وهذا هو \omterm{def:g11:drift-to-spreading:spreading-rate}{معدل التوسع الكلي} للظهرة في
آيسلندا. (3) وهو يتفق مع المعدل \qty{2}{cm} في السنة المذكور لشمال
الأطلسي.
\end{solution}

\begin{solution}{exo:g11:plate-motions:8}
(1) $v = \pi \times 0.651 \times 6371/180 \approx \qty{72.4}{km}$ في كل مليون
سنة، أي نحو \qty{7.2}{cm} في السنة. (2) $72.4 \times \sin\ang{30} \approx
\qty{36.2}{km}$ في كل مليون سنة، أي نحو \qty{3.6}{cm} في السنة. (3) عند قطب
الدوران، إن وقع على الصفيحة. فالسرعة تتغير من نقطة إلى
أخرى، من الصفر عند القطب إلى أكبر قيمة لها على بعد \ang{90} منه: فليس
للصفيحة كلها سرعة واحدة، بل سرعة عند كل نقطة (أو متوسط
على سطحها).
\end{solution}

\begin{solution}{exo:g11:plate-motions:9}
(1) بالملليمترات في السنة،
\[
  \vec v_{\mathrm{IN/EU}} = (40.6 - 24.0,\ 38.6 + 3.3) = (+16.6, +41.9).
\] والمقدار $\sqrt{16.6^2 + 41.9^2} \approx
\qty{45}{mm}$ في السنة؛ والزاوية $\beta$ مع الشمال تحقق $\tan\beta =
16.6/41.9 \approx 0.40$، إذن $\beta \approx \ang{22}$: فالهند تتحرك نحو
أوراسيا، في اتجاه يميل \ang{22} شرق الشمال. (2) \qty{45}{km} في كل مليون سنة، أي
\qty{450}{km} في 10 ملايين سنة.
\end{solution}

\begin{solution}{exo:g11:plate-motions:10}
$(2432 - 1058)/(27.7 - 10.3) = 1374/17.4 \approx \qty{79}{km}$ في كل مليون
سنة، أي نحو \qty{7.9}{cm} في السنة، وهذا أقل قليلًا من \qty{9.1}{cm} في السنة.
وفي كل من العمرين عدم يقين بنحو نصف مليون سنة؛ ففرق
العمرين، $17.4 \pm 1.0$ مليون سنة، يعطي سرعة بين نحو 75
و\qty{84}{km} في كل مليون سنة. ويبدو أن الصفيحة تحركت بين هذين البركانين
أبطأ قليلًا من المتوسط.
\end{solution}

\begin{solution}{exo:g11:plate-motions:11}
(1) $\vec v_{A/C} = \vec v_{A/B} + \vec v_{B/C}$: \qty{3}{cm} في السنة نحو الشرق مضافًا إليها
\qty{4}{cm} في السنة نحو الشمال. والسهمان يصنعان زاوية قائمة، وطول الضلع
الثالث للمثلث $\sqrt{3^2 + 4^2} = \qty{5}{cm}$ في السنة،
في اتجاه يميل \ang{37} شرق الشمال ($\tan\beta = 3/4$). (2) $\vec v_{C/A}
= -\vec v_{A/C}$: \qty{5}{cm} في السنة في الاتجاه المعاكس، أي في اتجاه يميل \ang{37} غرب
الجنوب.
\end{solution}

\begin{solution}{exo:g11:plate-motions:12}
(1) $54/0.375 = \qty{144}{km}$ في كل مليون سنة، أي \qty{14.4}{cm} في السنة. (2)
من $29/0.375 \approx 77$ إلى $79/0.375 \approx \qty{211}{km}$ في كل مليون
سنة: أي بين نحو 8 و\qty{21}{cm} في السنة. (3) قرب البقعة الساخنة لا تزيد
المسافة على حجم بركان أو على المنطقة التي تغذي فيها
البقعة الساخنة البراكين، ولا يزيد العمر على الزمن الذي يستغرقه بركان
لينمو: فكلاهما غير يقيني جدًا نسبيًا. أما بعيدًا عن
البقعة الساخنة فالخطأ بمقدار \qty{25}{km} على \qty{2400}{km}، أو بنصف مليون سنة
على 28 مليونًا، لا يتجاوز واحدًا أو اثنين في المئة.
\end{solution}

\begin{solution}{exo:g11:plate-motions:13}
(1) عشرة أضعاف التبعثر \qty{50}{mm}: $50/72 \approx 0.7$ سنة، أي نحو
ثمانية أشهر. (2) $50/5 = 10$ سنوات. (3) على مدى ملايين السنين تتحرك الصفائح
عشرات إلى مئات الكيلومترات، وهي مسافات تسجلها الشرائط المغناطيسية وكان
قياسها في البحر سهلًا؛ أما تحديد المواقع بدقة بضعة ملليمترات،
الذي يجعل الحركة مرئية في أشهر، فقد جاء بعد ذلك بزمن طويل.
\end{solution}

\begin{solution}{exo:g11:plate-motions:14}
(1) بالملليمترات في السنة،
\[
  \vec v_{\mathrm{SA/AF}} = (-4.9 - 27.1,\ 11.0 - 17.9) = (-32.0, -6.9).
\] والمقدار $\sqrt{32.0^2 + 6.9^2} \approx
\qty{32.7}{mm}$ في السنة، و$\tan\alpha = 6.9/32.0 \approx 0.22$ يعطي
اتجاهًا يميل \ang{12} جنوب الغرب. (2) \omterm{def:g11:drift-to-spreading:spreading-rate}{معدل التوسع الكلي} نحو \qty{3.3}{cm} في
السنة، ونصف المعدل نحو \qty{1.6}{cm} في السنة. (3) نصف المعدل في آخر
بضعة ملايين من السنين، \qty{1.6}{cm} في السنة، أصغر قليلًا من
المعدل \qty{1.9}{cm} في السنة الذي هو متوسط على 66 مليون سنة: فربما
انفتح جنوب الأطلسي في الآونة الأخيرة أبطأ قليلًا من المتوسط، أو ربما
يأتي الفرق من أوجه عدم اليقين في أعمار الرواسب. والقياسان كلاهما يتفقان
على أن الظهرة بطيئة.
\end{solution}

\begin{solution}{exo:g11:plate-motions:15}
(1) مستقيم يرتفع \qty{20}{mm} في السنة، ثم هبوط مفاجئ قدره
\qty{50}{mm} بعد خمس سنوات، ثم مستقيم يرتفع من جديد بالميل نفسه.
(2) تنتهي المحطة على بعد $10 \times 20 - 50 = \qty{150}{mm}$ شرق موضعها الأول:
$150/10 = \qty{15}{mm}$ في السنة، أي أصغر من الحقيقة بالربع. (3) يعطي كل نصف
\qty{20}{mm} في السنة. (4) القفزة حركة منفصلة مفاجئة سببها
الزلزال: فيقيس الجيوديسي مقدارها من المواقع التي تسبقها مباشرة
والتي تليها مباشرة، ويحذفها من السلسلة (أو يلائم الجزأين بالميل
نفسه مع درجة بينهما)، ويحتفظ بالقيمة \qty{20}{mm} في السنة سرعةً
للمحطة.
\end{solution}

\begin{solution}{pb:g11:plate-motions:1}
\textbf{1.} ماونا كيا (\num{0.375} مليون سنة) أحدثها، وسويكو
(الوسط) (\num{64.7} مليون سنة) أقدمها. وتزداد الأعمار مع
البعد عن كيلاويا.

\textbf{2.} $\qty{1}{km} = 10^6~\unit{mm}$ والمليون سنة $10^6$
سنة: فإن \qty{1}{km} في كل مليون سنة هو $10^6~\unit{mm}/10^6$ سنة،
أي \qty{1}{mm} في السنة.

\textbf{3.} $519/5.1 \approx \qty{102}{mm}$ في السنة.

\textbf{4.} نيكر $1058/10.3 \approx 102.7$؛ لايسان $1818/19.9 \approx
91.4$؛ ميدواي $2432/27.7 \approx 87.8$؛ كولاهان $3128/38.6 \approx 81.0$؛
أبوت $3280/38.7 \approx 84.8$ (بالملليمترات في السنة).

\textbf{5.} النقاط (54؛ 0.375)، (519؛ 5.1)، (780؛ 7.2)، (1058؛ 10.3)،
(1818؛ 19.9)، (2432؛ 27.7)، (3128؛ 38.6)، (3280؛ 38.7) تقع قرب مستقيم
يمر بالأصل، مع بعض التبعثر (فبالمقارنة بمستقيم \qty{91}{km} في كل
مليون سنة، تقع كاواي ونيهوا ونيكر تحته قليلًا، ولايسان عليه،
وكولاهان وأبوت فوقه قليلًا).

\textbf{6.} بُني كل بركان فوق البقعة الساخنة وتحرك مع
الصفيحة منذ ذلك الحين. وفي الزمن $t$ تحركت الصفيحة مسافة $v\,t$ في
اتجاه ثابت، فيقع البركان على بعد $d = v\,t$ من النقطة الواقعة الآن فوق
البقعة الساخنة: $t = d/v$.

\textbf{7.} $(102.7 + 91.4 + 87.8 + 81.0 + 84.8)/5 \approx \qty{89.5}{mm}$ في
السنة.

\textbf{8.} 89.5 مقابل \qty{91}{km} في كل مليون سنة: فهما يختلفان
بأقل من \qty{2}{\percent}.

\textbf{9.} $519/91 \approx 5.7$ مليون سنة، مقابل عمر قدره \num{5.1}
مليون سنة: وهما متقاربان، وكاواي أحدث قليلًا مما يتنبأ به
المستقيم المتوسط.

\textbf{10.} تزداد البراكين قِدمًا نحو الغرب الشمالي الغربي:
فالصفيحة تحملها في ذلك الاتجاه. فالصفيحة الهادئية تتحرك نحو
الغرب الشمالي الغربي فوق البقعة الساخنة.

\textbf{11.} $(4860 - 3493)/(64.7 - 42.4) = 1367/22.3 \approx \qty{61}{km}$
في كل مليون سنة، أي نحو \qty{6.1}{cm} في السنة.

\textbf{12.} $65/91 \approx 0.71$: فعلى امتداد سلسلة الإمبراطور تحركت الصفيحة
أبطأ بنحو \qty{30}{\percent} منها على امتداد سلسلة هاواي.

\textbf{13.} $3668/50 \approx \qty{73}{km}$ في كل مليون سنة، أي نحو
\qty{7.3}{cm} في السنة، وهذا أقل من \qty{91}{km} في كل مليون سنة. فالمستقيم المار بالأصل
يستعمل أعمار عام 1987، التي تضع نهاية سلسلة هاواي عند 39 إلى 43
مليون سنة؛ أما العمر الأحدث البالغ نحو 50 مليون سنة لبداية
الانحناء فيوزع المسافة نفسها على زمن أطول. والقيمتان صحيحتان كلتاهما إذا كانت الصفيحة
قد تحركت أبطأ في الجزء الأقدم من سلسلة هاواي وأسرع بعد ذلك:
فالسرعة لم تكن ثابتة.

\textbf{14.} تغيرًا في اتجاه حركة الصفيحة الهادئية، قبل نحو
50 مليون سنة.

\textbf{15.} اتجاه السلسلة هو الاتجاه الذي تحركت فيه
الصفيحة؛ وتباعد الأعمار على امتدادها هو سرعتها. ويمكن للصفيحة أن تتسارع أو
تتباطأ دون أن تنعطف، أو أن تنعطف دون أن تغيّر سرعتها، فالمعلومتان
مستقلتان. وتغيرات السرعة دون تغير في الاتجاه على امتداد سلسلة
الإمبراطور تبيّن أن سجل سرعتها والانحناء حدثان منفصلان: فالانحناء
لم يصنعه تباطؤ البقعة الساخنة.

\textbf{16.} $\sqrt{63.2^2 + 34.5^2} \approx \qty{72.0}{mm}$ في السنة،
و$\tan\alpha = 34.5/63.2 \approx 0.55$ يعطي اتجاهًا يميل \ang{29} شمال الغرب.

\textbf{17.} الغرب الشمالي الغربي يميل \ang{22.5} شمال الغرب، والشمال الغربي
\ang{45}: فاتجاه المحطة، \ang{29} شمال الغرب، يقع بين
الاثنين، قريبًا من اتجاه السلسلة. فهو متسق معها.

\textbf{18.} تقاس سرعة المحطة بالنسبة إلى إطار لا يدور فيه
\omterm{def:g11:lithosphere-asthenosphere:lithosphere}{الغلاف الصخري} ككل، في حين تعطي السلسلة
الحركة بالنسبة إلى البقعة الساخنة: فإذا تحرك الإطاران أحدهما بالنسبة إلى
الآخر، اختلفت السرعتان. ثم إن المحطة تقيس آخر 30 سنة،
والسلسلة آخر 40 مليون سنة، التي قد تكون السرعة تغيرت خلالها.
ويترك تبعثر الأعمار أيضًا عدم يقين بضعة ملليمترات في السنة
في سرعة السلسلة.

\textbf{19.} $100/91 \approx 1.1$ مليون سنة.

\textbf{20.} \textbf{تحركت الصفيحة الهادئية فوق بقعة هاواي الساخنة بنحو
\qty{9}{cm} في السنة} خلال آخر 40 مليون سنة، نحو
الغرب الشمالي الغربي؛ وقبل الانحناء كانت تتحرك أبطأ، وفي اتجاه
آخر.
\end{solution}
